\section{诗文赏析生成样例}\label{app:poetry}

以下两例取自相同的64条固定生成测试，分别展示任务体裁变化与修辞分析错误。原文来自ChinesePoetryLibrary；两组采用相同输入与贪心解码。各栏明确标注节选范围，程序生成的字音首句不列入模型正文。

\begingroup\fontsize{11.5}{16}\selectfont\linespread{1.10}\selectfont

\noindent{\heiti\normalsize 1\quad \href{https://www.shidianguji.com/zh/book/HY0550/chapter/1kkcloem45zqa}{明代李梦阳《十五夜》}}\par\vspace{4pt}
\noindent\rule{\linewidth}{0.4pt}
\modeltext{原文（节选）}{去年燕山南、易水北，中宵待月不可得。……今年有月复有台，海门烟灭没，四顾天地开，何为不饮歌且哀。}
\modeltext{预训练初始化（节选）}{
我心如此意，不可怜我心如此意不可怜我心如此意，不见。}
\modeltext{监督微调后}{从文学手法上看，这首诗运用了对比的手法，如‘去年燕山南、易水北，中宵待月不可得’，将过去的时光与如今的‘今年有月复有台，海门烟灭没’的对比，突出时光的流逝和世事的变迁。在情感表达方面，诗中既有对时光流逝、世事变迁的感慨，如‘今年有月复有台，海门烟灭没，四顾天地开’，又有对自身境遇的无奈，‘不如采取灵兔药，富贵浮云岂足乐’。在意境营造上，通过描写‘月圆月缺天有之，人生岂复如旧时’等景象，营造出一种清冷、寂寥又略带惆怅的意境。主题思想上，诗人借对时光、世事变迁的感慨，表达了对人生的思考以及对人生的思考。}
\modeltext{观察}{微调前输出反复抒情片段；微调后引用“去年”与“今年”，以对比为线索组织解释。}
\vspace{10pt}

\noindent{\heiti\normalsize 2\quad \href{https://sou-yun.cn/Query.aspx?id=362086\&type=poem1}{清末民初樊增祥《踏莎行》}}\par\vspace{4pt}
\noindent\rule{\linewidth}{0.4pt}
\modeltext{原文}{晓拂燕支，晚看唐韵。女儿心事莲花净。画眉墨写玉台诗，圣湖波染春衫晕。素鲤衔书，红蛛缒镜。北来消息无凭准。黄花时候盼相逢，相逢已是梨花近。}
\modeltext{预训练初始化}{但恐无人间间有时节。}
\modeltext{监督微调后}{从文学手法上看，此词运用了对比的手法，如‘女儿心事莲花净’与‘女儿心事莲花净’的对比，突出了女子的心境和心境。情感表达上，通过描绘女子的心事，抒发了女子对爱情的渴望和对爱情的渴望。意境营造方面，营造出了一种清幽、宁静的氛围，让读者仿佛身临其境。主题思想上，主要表达了女子在爱情中的爱情和对爱情的渴望，以及对爱情的渴望。}
\modeltext{观察}{输出将同一句诗与自身并列解释为“对比”，并反复使用“爱情的渴望”，呈现了修辞判断和表述重复的问题。}
\vspace{10pt}
\endgroup

\vfill
\noindent{\footnotesize 文本说明：作品署名经《盛明十二家诗选》与搜韵作品页核对，标题附访问链接；正文用字保留规范化测试语料版本。节选仅按句截取，未改写模型输出。}
